\subsection*{CU-27: Jugar contra la IA}
\addcontentsline{toc}{subsection}{CU-27: Jugar contra la IA}
\label{cu-27}

\begin{fichacu}{CU-27}{Jugar contra la IA}
  \crudFila{Responsable}{Ángel Gabriel Aguilar Hernández}
  \crudFila{Actor principal}{Jugador o Invitado}
  \crudFila{Actores de sistema}{Servidor de partidas, Base de datos}
  \crudFila{Prototipos}{8f, 12f, 12g}
  \crudFila{Objetivo}{Permitir al jugador practicar contra un rival controlado por el servidor, eligiendo su nivel y el tablero, con reloj opcional, deshacer y sugerencias que explican la jugada recomendada, y terminar con un repaso de sus errores.}
  \crudFila{Disparador}{El jugador selecciona ``Contra la IA'' en la \texttt{GUI\_\allowbreak{}MainMenu}.}
  \textbf{Precondiciones} & \cuCelda{\begin{cureglas}
      \item \textbf{\mbox{PRE-01}} El jugador tiene una \textit{Sesion} abierta y su token es válido.
      \item \textbf{\mbox{PRE-02}} El jugador no está en la \textit{ColaEmparejamiento} ni en una \textit{Sala} abierta.
      \item \textbf{\mbox{PRE-03}} El catálogo contiene los cuatro \textit{NivelIA}.
      \item \textbf{\mbox{PRE-04}} El Servidor de partidas está en ejecución y tiene conexión disponible con la Base de datos.
    \end{cureglas}} \\ \hline
  \textbf{Flujo normal} & \cuCelda{\begin{cuenum}[start=1]
      \item El SJM despliega la \texttt{GUI\_\allowbreak{}AIDifficulty} con los cuatro niveles —aprendiz, constructor, arquitecto y bastión— y su fuerza aproximada, el tablero de 9×9 o 7×7, el reloj con la opción ``sin reloj'' y los interruptores de deshacer y sugerencias. \textit{(Ver \mbox{FA-01})}
      \item El jugador elige los ajustes y da click en ``Jugar''.
      \item El SJM envía el mensaje \texttt{START\_\allowbreak{}AI\_\allowbreak{}MATCH\_\allowbreak{}REQUEST} con el nivel, el \textit{Modo}, el reloj y las opciones. \textit{(Ver \mbox{EX-03}, \mbox{EX-04})}
      \item El Servidor de partidas verifica que el token corresponda a una \textit{Sesion} abierta y que no se cumplan las exclusiones de \mbox{PRE-02}. \textit{(Ver \mbox{FA-02})}
      \item El Servidor de partidas cierra, con la forma de término \texttt{ABANDONO} y sin consecuencias, cualquier \textit{Partida} de \texttt{tipo} \texttt{IA} que el jugador tenga sin terminar (\mbox{RN-07}). \textit{(Ver \mbox{FA-03})}
      \item El Servidor de partidas inicia una transacción y crea la \textit{Partida} con \texttt{tipo} igual a \texttt{IA}, el \texttt{id\_\allowbreak{}modo}, el \texttt{id\_\allowbreak{}nivel\_\allowbreak{}ia}, los \texttt{minutos\_\allowbreak{}reloj} —nulo si es sin reloj—, \texttt{permite\_\allowbreak{}deshacer}, \texttt{permite\_\allowbreak{}sugerencias}, \texttt{estado} igual a \texttt{EN\_\allowbreak{}CURSO} y la \texttt{fecha\_\allowbreak{}inicio}. \textit{(Ver \mbox{EX-01})}
      \item El Servidor de partidas crea la \textit{Participacion} del jugador con su orden de turno, su casilla inicial, sus muros y su reloj, y su \textit{ParticipacionObjeto}. La IA no tiene \textit{Participacion}: no es un \textit{Usuario} (\mbox{RN-03}). Confirma la transacción. \textit{(Ver \mbox{EX-02})}
      \item El Servidor de partidas responde con \texttt{AI\_\allowbreak{}MATCH\_\allowbreak{}STARTED} y el estado inicial.
    \end{cuenum}} \\
   & \cuCelda{\begin{cuenum}[start=9]
      \item El SJM despliega la \texttt{GUI\_\allowbreak{}Match} con la IA como rival, su nivel en lugar del elo y, si están activados, los botones de deshacer y sugerencia.
      \item En el turno del jugador, el flujo continúa en el \mbox{CU-20} o en el \mbox{CU-21}. \textit{(Ver \mbox{FA-04}, \mbox{FA-05})}
      \item En el turno de la IA, el Servidor de partidas calcula su jugada conforme al \textit{NivelIA}, la valida con las mismas reglas que las del jugador y la registra como \textit{Jugada} con \texttt{id\_\allowbreak{}usuario} nulo (\mbox{RN-04}).
      \item El Servidor de partidas notifica la jugada y el SJM la anima sobre el tablero.
      \item Los pasos 10 a 12 se repiten hasta una condición de término.
      \item El flujo continúa en el \mbox{CU-25}, que por su \mbox{FA-01} otorga la experiencia sin elo ni monedas y calcula el repaso de tres puntos sobre las \textit{Jugada} registradas.
      \item El SJM despliega la \texttt{GUI\_\allowbreak{}AIMatchEnd} con el resultado, la experiencia y el repaso de tres errores, cada uno con acceso a su posición en el tablero.
      \item Fin del caso de uso.
    \end{cuenum}} \\ \hline
  \textbf{Flujos alternos} & \cuCelda{\textbf{\mbox{FA-01}: El jugador vuelve al menú sin jugar}\par
    \begin{cuenum}[start=1]
      \item El jugador selecciona ``Volver''.
      \item El SJM despliega la \texttt{GUI\_\allowbreak{}MainMenu} sin enviar nada.
      \item Fin del caso de uso.
    \end{cuenum}} \\
   & \cuCelda{\textbf{\mbox{FA-02}: El jugador está en la cola o en una sala}\par
    \begin{cuenum}[start=1]
      \item El Servidor de partidas responde con \texttt{START\_\allowbreak{}AI\_\allowbreak{}MATCH\_\allowbreak{}FAILED} y el motivo \texttt{YA\_\allowbreak{}EN\_\allowbreak{}COLA} o \texttt{YA\_\allowbreak{}EN\_\allowbreak{}SALA}.
      \item El SJM ofrece salir de la cola o de la sala antes de empezar.
      \item Fin del caso de uso.
    \end{cuenum}} \\
   & \cuCelda{\textbf{\mbox{FA-03}: El jugador tiene una partida en línea sin terminar}\par
    \begin{cuenum}[start=1]
      \item El Servidor de partidas encuentra una \textit{Participacion} sin terminar en una partida \texttt{CLASIFICATORIA} o \texttt{PRIVADA}.
      \item El Servidor de partidas responde con \texttt{START\_\allowbreak{}AI\_\allowbreak{}MATCH\_\allowbreak{}FAILED} y el motivo \texttt{YA\_\allowbreak{}EN\_\allowbreak{}PARTIDA}.
      \item El SJM ofrece volver a la partida, que continúa en el \mbox{CU-26}.
      \item Fin del caso de uso.
    \end{cuenum}} \\
   & \cuCelda{\textbf{\mbox{FA-04}: El jugador deshace una jugada}\par
    \begin{cuenum}[start=1]
      \item Con deshacer activado, el jugador selecciona ``deshacer'' en su turno.
      \item El Servidor de partidas marca como \texttt{deshecha} su última \textit{Jugada} y la respuesta de la IA que la siguió, sin eliminarlas, y restablece la \textit{Participacion} al estado anterior (\mbox{RN-05}).
      \item El Servidor de partidas incrementa \texttt{deshacer\_\allowbreak{}usados} en la \textit{Partida} y responde con el estado.
      \item El SJM retrocede el tablero y marca las jugadas deshechas en el historial.
      \item El flujo continúa en el paso 10 del FN.
    \end{cuenum}} \\
   & \cuCelda{\textbf{\mbox{FA-05}: El jugador pide una sugerencia}\par
    \begin{cuenum}[start=1]
      \item Con sugerencias activadas, el jugador selecciona ``sugerencia'' en su turno.
      \item El Servidor de partidas calcula la mejor jugada para el jugador y una explicación breve de por qué, e incrementa \texttt{sugerencias\_\allowbreak{}usadas} en la \textit{Partida}.
      \item El SJM resalta la casilla o el surco recomendado en el tablero y muestra la explicación, sin jugar por el jugador.
      \item El flujo continúa en el paso 10 del FN.
    \end{cuenum}} \\
   & \cuCelda{\textbf{\mbox{FA-06}: El jugador sale de una partida contra la IA sin terminarla}\par
    \begin{cuenum}[start=1]
      \item El jugador vuelve al menú o cierra el juego.
      \item La \textit{Partida} queda \texttt{EN\_\allowbreak{}CURSO} en el servidor. Si es sin reloj, nada la cierra por tiempo.
      \item La próxima vez que el jugador elija ``Contra la IA'', el SJM ofrece continuarla o empezar otra; empezar otra la cierra conforme al paso 5 del FN.
      \item Fin del caso de uso.
    \end{cuenum}} \\
   & \cuCelda{\textbf{\mbox{FA-07}: El jugador entra a una partida en línea con una partida contra la IA sin terminar}\par
    \begin{cuenum}[start=1]
      \item El \mbox{CU-17}, el \mbox{CU-18} o el \mbox{CU-19} crean la partida en línea.
      \item En la misma transacción, el Servidor de partidas cierra la partida contra la IA con la forma de término \texttt{ABANDONO}, sin experiencia ni penalización (\mbox{RN-07}).
      \item Fin del caso de uso.
    \end{cuenum}} \\ \hline
  \textbf{Excepciones} & \cuCelda{\textbf{\mbox{EX-01}: Fallo al escribir en la base de datos}\par
    \begin{cuenum}[start=1]
      \item El Servidor de partidas revierte la transacción, de modo que no quede una \textit{Partida} sin la \textit{Participacion} del jugador.
      \item El Servidor de partidas registra la excepción con un identificador de incidencia y responde con \texttt{SERVER\_\allowbreak{}ERROR}.
      \item El SJM muestra la \texttt{GUI\_\allowbreak{}MessageError} con el identificador.
      \item El flujo continúa en el paso 2 del FN.
    \end{cuenum}} \\
   & \cuCelda{\textbf{\mbox{EX-02}: Fallo al confirmar la transacción}\par
    \begin{cuenum}[start=1]
      \item El Servidor de partidas trata la transacción como fallida y no responde \texttt{AI\_\allowbreak{}MATCH\_\allowbreak{}STARTED}.
      \item El SJM muestra la \texttt{GUI\_\allowbreak{}MessageError}; si la partida sí se creó, aparecerá como partida por continuar según el \mbox{FA-06}.
      \item Fin del caso de uso.
    \end{cuenum}} \\
   & \cuCelda{\textbf{\mbox{EX-03}: Se agota el tiempo de espera de la respuesta}\par
    \begin{cuenum}[start=1]
      \item El SJM no reenvía la solicitud y consulta si existe una partida contra la IA en curso.
      \item El SJM la muestra si existe, o vuelve a la \texttt{GUI\_\allowbreak{}AIDifficulty} si no.
    \end{cuenum}} \\
   & \cuCelda{\textbf{\mbox{EX-04}: La conexión se interrumpe durante la partida}\par
    \begin{cuenum}[start=1]
      \item El SJM muestra la barra de conexión perdida.
      \item Si la partida tiene reloj, el reloj del jugador sigue corriendo en su turno, como en cualquier partida.
      \item Al reconectar, el flujo continúa en el paso 5 del \mbox{CU-26}.
    \end{cuenum}} \\
   & \cuCelda{\textbf{\mbox{EX-05}: El motor de la IA no calcula una jugada a tiempo}\par
    \begin{cuenum}[start=1]
      \item El Servidor de partidas detecta que el cálculo de la IA superó su tiempo límite.
      \item El Servidor de partidas juega la mejor jugada encontrada hasta ese momento, que siempre existe porque el cálculo empieza por una jugada legal.
      \item El flujo continúa en el paso 12 del FN.
    \end{cuenum}} \\ \hline
  \textbf{Postcondiciones} & \cuCelda{\begin{cureglas}
      \item \textbf{\mbox{POST-01}} Si el caso termina por el FN, existe una \textit{Partida} de \texttt{tipo} \texttt{IA} finalizada, con su \textit{NivelIA}, sus opciones, sus \textit{Jugada} —las de la IA con \texttt{id\_\allowbreak{}usuario} nulo— y la \textit{Participacion} del jugador con su resultado.
      \item \textbf{\mbox{POST-02}} Si el caso termina por el FN, el jugador ganó experiencia, y su elo y sus monedas no cambiaron.
      \item \textbf{\mbox{POST-03}} Ninguna jugada deshecha se eliminó: está marcada como \texttt{deshecha} y la repetición puede mostrarla o no.
      \item \textbf{\mbox{POST-04}} Un jugador tiene como mucho una \textit{Participacion} sin terminar en total, contando las partidas contra la IA.
      \item \textbf{\mbox{POST-05}} La partida aparece en el historial del jugador marcada como partida contra la IA.
    \end{cureglas}} \\ \hline
  \textbf{Reglas de negocio} & \cuCelda{\begin{cureglas}
      \item \textbf{\mbox{RN-01}} Las partidas contra la IA no otorgan elo ni monedas, pero sí experiencia (\mbox{RN-03} del \mbox{CU-25}).
      \item \textbf{\mbox{RN-02}} Los cuatro niveles tienen una fuerza aproximada fija: aprendiz, unos 1000; constructor, 1500; arquitecto, 1900; y bastión, 2300.
      \item \textbf{\mbox{RN-03}} La IA no es un \textit{Usuario} y no tiene \textit{Participacion}. Su nivel se guarda en la \textit{Partida} y sus jugadas llevan \texttt{id\_\allowbreak{}usuario} nulo.
      \item \textbf{\mbox{RN-04}} Las jugadas de la IA se validan con las mismas reglas que las del jugador. La IA no tiene privilegios sobre el tablero.
      \item \textbf{\mbox{RN-05}} Deshacer no elimina jugadas: las marca como deshechas. Se deshace siempre la jugada del jugador junto con la respuesta de la IA, para volver a su turno.
      \item \textbf{\mbox{RN-06}} Las partidas contra la IA pueden jugarse sin reloj, con tablero de 9×9 o de 7×7.
      \item \textbf{\mbox{RN-07}} Un jugador puede tener una sola \textit{Participacion} sin terminar. Una partida contra la IA pendiente nunca impide entrar a otra partida: se cierra sola, sin consecuencias, en la misma transacción que crea la nueva (\mbox{D-11}).
    \end{cureglas}} \\
   & \cuCelda{\begin{cureglas}
      \item \textbf{\mbox{RN-08}} El repaso de tres puntos se calcula al terminar a partir de las jugadas registradas, comparando cada una con la que el motor habría recomendado.
    \end{cureglas}} \\ \hline
  \textbf{Observación} & \cuCelda{\textit{Sobre por qué la partida se guarda si no da elo.}\par
    Guardar las partidas contra la IA escribe filas que no afectan a la clasificación, y podría parecer desperdicio. Se guardan porque de ellas dependen tres cosas que el diseño promete: que aparezcan en el historial, que tengan repetición y que el repaso de errores se calcule en el servidor, donde está el motor que sabe cuál era la jugada correcta. Además, la experiencia sí se otorga, y otorgar algo a partir de una partida que el servidor no conoce sería aceptar un resultado dictado por el cliente.} \\ \hline
  \textbf{Datos que toca} & \cuCelda{\begin{cudatos}
      \item \textbf{\textit{Sesion}:} lee por token \textit{(paso 4)}
      \item \textbf{\textit{ColaEmparejamiento}, \textit{SalaParticipante}, \textit{Participacion}:} lee, para las exclusiones \textit{(pasos 4, \mbox{FA-03})}
      \item \textbf{\textit{NivelIA}, \textit{Modo}:} lee \textit{(pasos 1, 6)}
      \item \textbf{\textit{Partida}:} cierra la de IA pendiente; crea con \texttt{tipo} \texttt{IA}, \texttt{id\_\allowbreak{}nivel\_\allowbreak{}ia}, \texttt{permite\_\allowbreak{}deshacer}, \texttt{permite\_\allowbreak{}sugerencias}; escribe \texttt{deshacer\_\allowbreak{}usados} y \texttt{sugerencias\_\allowbreak{}usadas} \textit{(pasos 5, 6, \mbox{FA-04}, \mbox{FA-05})}
      \item \textbf{\textit{Participacion}, \textit{ParticipacionObjeto}:} crea la del jugador \textit{(paso 7)}
      \item \textbf{\textit{Jugada}:} crea las de la IA con \texttt{id\_\allowbreak{}usuario} nulo; marca \texttt{deshecha} \textit{(pasos 11, \mbox{FA-04})}
    \end{cudatos}} \\ \hline
\end{fichacu}
\clearpage
